\section{The Traffic, in Full}
\label{sec:ch07-the-traffic-in-full}
\index{wire traffic|(}%

Six HTTP messages answer the chapter's query, and
figure~\ref{fig:ch07-message-sequence} is all of them in order. This section is
the two that go to the subgraphs, printed as the subgraphs recorded them.

\begin{figure}[htbp]
  \centering
  % Every HTTP message exchanged for
% { products { title price reviews { rating body } } }, in the order the
% subgraph consoles recorded them. Referenced from chapter 07. Bare
% tikzpicture; the figure environment, caption and label live at the call site.
\begin{tikzpicture}[
    font=\small,
    party/.style={draw, rounded corners=2pt, minimum width=24mm,
                  minimum height=8mm, align=center},
    life/.style={gray, dashed},
    msg/.style={-{Stealth[length=2mm]}},
    back/.style={-{Stealth[length=2mm]}, gray},
    lbl/.style={font=\scriptsize, align=center, fill=white, inner sep=1pt}
  ]

  \node[party] (client) at (0,0)    {Client};
  \node[party] (router) at (3.6,0)  {Router};
  \node[party] (cat)    at (7.9,0)  {Catalog\\\scriptsize :5201};
  \node[party] (rev)    at (12.2,0) {Reviews\\\scriptsize :5202};

  \foreach \x in {0, 3.6, 7.9, 12.2} {
    \draw[life] (\x,-0.5) -- (\x,-6.4);
  }

  % 1. the client's document
  \draw[msg] (0,-1.0) -- (3.6,-1.0);
  \node[lbl, anchor=south] at (1.8,-0.95) {\code{POST /graphql}};

  % 2 and 3. fetch 0
  \draw[msg] (3.6,-2.0) -- (7.9,-2.0);
  \node[lbl, anchor=south] at (5.75,-1.95)
    {\code{\{products \{title price}\\\code{\_\_typename id\}\}}};
  \draw[back] (7.9,-3.0) -- (3.6,-3.0);
  \node[lbl, anchor=south] at (5.75,-2.95) {3 products, with ids};

  % 4 and 5. fetch 1
  \draw[msg] (3.6,-4.0) -- (12.2,-4.0);
  \node[lbl, anchor=south] at (7.9,-3.95)
    {\code{\_entities}, 3 representations in one call};
  \draw[back] (12.2,-5.0) -- (3.6,-5.0);
  \node[lbl, anchor=south] at (7.9,-4.95) {3 review lists, one of them empty};

  % 6. the assembled answer
  \draw[back] (3.6,-6.0) -- (0,-6.0);
  \node[lbl, anchor=south] at (1.8,-5.95) {one response};

\end{tikzpicture}

  \caption{Six messages, four of them between the router and the subgraphs. The
  two fetches are strictly sequential because the second one needs identifiers
  the first one returns. Nothing here is proprietary: both subgraph calls are
  ordinary GraphQL over HTTP.}
  \label{fig:ch07-message-sequence}
\end{figure}

\index{wire traffic!capturing it}%
Before the listings, a word on where they come from, because it changes how much
weight they can carry. Both subgraphs call ASP.NET Core's own HTTP logging
middleware with request and response bodies turned on. There is no proxy in the
path and no packet capture. What follows is what the server process received and
what it sent, read from inside the server.

Two settings make that work and neither is obvious. A body is logged only if its
media type is on a list, and GraphQL over HTTP answers with
\code{application/graphql-response+json}, which ASP.NET Core does not know
about, so it has to be added by hand. And headers outside the middleware's
default allow-list are printed as \code{[Redacted]} rather than dropped, which
turns out to matter later in this section.

\subsection*{What Catalog received}

\begin{minted}[fontsize=\footnotesize]{text}
Method: POST
Path: /graphql
Accept: application/json
Host: host.docker.internal:5201
User-Agent: Go-http-client/1.1
Accept-Encoding: gzip,deflate
Content-Type: application/json
traceparent: 00-14c8c3610596717bc4c5d5d57ec4c99a-b1ecebfd8985308e-01
Content-Length: 50

RequestBody: {"query":"{products {title price __typename id}}"}
\end{minted}

Fifty bytes. An ordinary \code{POST} carrying a \code{query} key, with no
variables and no header a GraphQL service would have to be taught about. A
subgraph needs no router-specific code to be talked to, which has a consequence
worth holding on to: whatever you could do to a HotChocolate service in
chapters~\ref{ch:hotchocolate-quickly}
to~\ref{ch:schema-design-that-survives-change}, you can still do to it once it
is behind a router.

And what it sent back:

\begin{minted}[fontsize=\footnotesize]{text}
ResponseBody: {"data":{"products":[
  {"title":"Oak dining table","price":749.00,"__typename":"Product","id":"1"},
  {"title":"Linen armchair","price":429.00,"__typename":"Product","id":"2"},
  {"title":"Brass floor lamp","price":189.00,"__typename":"Product","id":"3"}]}}
\end{minted}

The \code{\_\_typename} and \code{id} the router added in
section~\ref{sec:ch07-the-plan} come back in the payload, get stripped out
before the client sees anything, and exist only to make the next message
possible.

\subsection*{What Reviews received}

\begin{minted}[fontsize=\footnotesize]{text}
traceparent: 00-14c8c3610596717bc4c5d5d57ec4c99a-ad519fea48de7013-01
Content-Length: 278

RequestBody: {"variables":{"representations":[
  {"__typename":"Product","id":"1"},
  {"__typename":"Product","id":"2"},
  {"__typename":"Product","id":"3"}]},
  "query":"query($representations: [_Any!]!){_entities(representations:
  $representations){... on Product {__typename reviews {rating body}}}}"}
\end{minted}

\index{wire traffic!batching}%
One call for three products. This is the fact the whole arrangement stands on,
and the first thing I would check in a federated graph I inherited. Three
separate calls would produce an identical answer to the client and turn a
two-request query into a four-request query, and nothing in the response would
tell you it had happened. The verification script in the companion repository
asserts this body character for character for exactly that reason: an assertion
on the answer would pass either way.

\subsection*{One trace, two spans}

Both subgraph requests carry a \code{traceparent}, and the two are not the same:

\begin{minted}[fontsize=\footnotesize]{text}
catalog:  00-14c8c3610596717bc4c5d5d57ec4c99a-b1ecebfd8985308e-01
reviews:  00-14c8c3610596717bc4c5d5d57ec4c99a-ad519fea48de7013-01
\end{minted}

\index{tracing!W3C trace context}%
Same trace identifier, different span identifiers. The router generated both.
That is W3C trace context doing the one job that makes a federated request
debuggable at all: without it a slow query is three unrelated log entries in
three services, and with it they are one tree.
Chapter~\ref{ch:observability} builds the rest of that picture. What is worth
knowing here is that the router does this whether or not anything is collecting
traces, so the identifiers are in your subgraph logs already.

\subsection*{Nothing else of the client's request survives}

I sent the router a request carrying three headers a real client might send:

\begin{minted}[fontsize=\footnotesize]{text}
graphql-client-name: chapter-07
Authorization: Bearer not-a-real-token
X-Custom-Thing: hello
\end{minted}

None of the three reached the subgraph. Not redacted, not renamed: absent.

\index{router!header propagation}%
That is a measurement rather than a reading, and it is worth explaining how it
can be one. Recall that the logging middleware prints an unknown header as
\code{[Redacted]} rather than dropping it. Sent straight to the subgraph, the
same three headers log like this:

\begin{minted}[fontsize=\footnotesize]{text}
Authorization: [Redacted]
Content-Type: application/json
graphql-client-name: chapter-07
X-Custom-Thing: [Redacted]
\end{minted}

Through the router, none of the three lines appears at all. A header the router
forwarded but that the log declined to print would look different from a header
the router never sent, which is what makes the absence evidence.

The documentation agrees: \enquote{By default, no headers are forwarded for
security reasons}~\autocite{cosmo2026headers}. Turning one on is a router
configuration rule, not a code change:

\begin{minted}{yaml}
headers:
  all:
    request:
      - op: "propagate"
        named: X-Test-Header
\end{minted}

\index{authorisation!at the boundary}%
I think this default is right, and it is the first place a federated graph will
surprise a team that has been passing an \code{Authorization} header around a
monolith without thinking about it. The router is a trust boundary now. Deciding
what crosses it is a design decision with a security consequence, and
chapter~\ref{ch:identity-and-authorization} is the whole of that decision rather
than the one-line configuration that implements it.

\subsection*{Content negotiation, briefly}

The router asks for \code{application/json} and gets it, with a charset
attached. It does not ask for the newer
\code{application/graphql-response+json}, which is what HotChocolate would
answer a client that did. Both are correct under the GraphQL over HTTP
specification and they differ in how status codes are chosen, which is a real
distinction with real consequences and belongs to
chapter~\ref{ch:the-wire-completely}.

\index{wire traffic|)}%
